\section{Results}
\label{sec:results}

\subsection{Study cohort}

The cohort comprised 48,354 admissions among 29,577 unique adult patients
across six Israeli general hospitals, with center-specific sample sizes ranging
from 4,478 to 11,024 admissions (Table~\ref{tab:cohort}). Only 21.5\% of
admissions meeting a hyperglycemia criterion carried a documented diabetes
diagnosis, indicating that a substantial proportion of the cohort represented
clinically significant inpatient hyperglycemia without a recorded diagnosis.

\begin{table}[t]
\centering
\caption{Characteristics of the multicenter study cohort.}
\label{tab:cohort}
\small
\begin{tabularx}{\textwidth}{
>{\raggedright\arraybackslash}X
>{\raggedleft\arraybackslash}p{0.34\textwidth}
}
\toprule
Characteristic & Value \\
\midrule
Hospitals & 6 \\
Admissions & 48,354 \\
Unique patients & 29,577 \\
Age, mean (median), y & 71.9 (73.3) \\
Male sex & 54.3\% \\
Length of stay, median, d & 5.55 \\
30-day mortality & 9.58\% \\
\midrule
Documented diabetes diagnosis & 35.3\% \\
Recurrent hyperglycemia & 74.7\% \\
Extreme glucose ($\geq$250 mg/dL) & 60.2\% \\
Hyperglycemia among admissions with diabetes & 50.3\% \\
Diabetes diagnosis among admissions with hyperglycemia & 21.5\% \\
\midrule
Pre-intervention cohort &
\makecell[r]{21,605 admissions\\17,002 patients} \\
Post-intervention cohort &
\makecell[r]{16,182 admissions\\11,201 patients} \\
\bottomrule
\end{tabularx}
\end{table}

\subsection{Adherence and its change following national implementation}
\label{sec:adherence-results}

Adherence varied sharply across recommendations, and the split followed the
kind of action required rather than its clinical weight
(Table~\ref{tab:adherence}). Single discrete actions on admission were
near-universal: SGLT2 discontinuation was complete (1.00) and creatinine
measurement almost so (0.98). Recommendations calling for repeated insulin
decisions across the stay were not: insulin following hyperglycemia scored 0.15, the
two routine basal dosing components 0.21 and 0.23, and bolus reduction after
hypoglycemia 0.28---all with a median of zero, meaning the expected action
was absent altogether in most eligible admissions. Monitoring and admission
medication management fell between these extremes (0.54--0.88). Adherence was
largely independent of whether diabetes had been formally documented (0.669 vs
0.661 among admissions with hyperglycemia in the first 48 hours), a difference
reaching significance only by virtue of the sample size.

\begin{table}[t]
\centering
\caption{Adherence to each assessed inpatient diabetes guideline
recommendation over the whole study period, and its change following
national implementation. Scores range from 0 (nonadherence) to 1 (complete
adherence). $\Delta$ is the absolute difference in mean adherence between the
post- and pre-intervention periods, with unadjusted 95\% CIs; $P$ values were
adjusted across component-level comparisons using the Benjamini--Hochberg
false discovery rate procedure.}
\label{tab:adherence}
\small
\setlength{\tabcolsep}{4pt}
\begin{tabularx}{\textwidth}{
>{\raggedright\arraybackslash}X
rrrrlr
}
\toprule
Recommendation & Mean & Median & Before & After & $\Delta$ (95\% CI) & $P$ \\
\midrule
Admission glucose monitoring
  & 0.85 & 1.00 & 0.78 & 0.93 & $+$0.15 (0.145 to 0.154) & $<.001$ \\
Routine glucose monitoring
  & 0.88 & 0.98 & 0.83 & 0.94 & $+$0.11 (0.105 to 0.113) & $<.001$ \\
Glucose monitoring during corticosteroid therapy
  & 0.54 & 0.50 & 0.51 & 0.58 & $+$0.07 (0.062 to 0.087) & $<.001$ \\
BMI measurement on admission
  & 0.87 & 1.00 & 0.86 & 0.89 & $+$0.03 (0.015 to 0.029) & $<.001$ \\
Creatinine measurement on admission
  & 0.98 & 1.00 & 0.98 & 0.99 & $+$0.01 (0.008 to 0.014) & $<.001$ \\
Antidiabetic discontinuation on admission
  & 0.66 & 0.86 & 0.62 & 0.75 & $+$0.13 (0.075 to 0.192) & $<.001$ \\
SGLT2 discontinuation on admission
  & 1.00 & 1.00 & 1.00 & 1.00 & 0.00 (--) & -- \\
Basal insulin continuation on admission
  & 0.66 & 1.00 & 0.63 & 0.67 & $+$0.04 (0.002 to 0.089) & .054 \\
Insulin administration after hyperglycemia
  & 0.15 & 0.00 & 0.14 & 0.15 & $+$0.01 (0.009 to 0.020) & $<.001$ \\
Antidiabetic discontinuation after hypoglycemia
  & 0.89 & 1.00 & 0.91 & 0.89 & $-$0.02 ($-$0.197 to 0.158) & .818 \\
Basal dose reduction after hypoglycemia
  & 0.87 & 1.00 & 0.87 & 0.88 & $+$0.01 ($-$0.015 to 0.034) & .497 \\
Bolus dose reduction after hypoglycemia
  & 0.28 & 0.00 & 0.33 & 0.26 & $-$0.07 ($-$0.158 to 0.029) & .210 \\
Routine basal dosage (frequency)
  & 0.23 & 0.00 & 0.21 & 0.24 & $+$0.03 (0.027 to 0.043) & $<.001$ \\
Routine basal dosage (dose)
  & 0.21 & 0.00 & 0.19 & 0.22 & $+$0.03 (0.021 to 0.036) & $<.001$ \\
\bottomrule
\end{tabularx}
\end{table}

Overall adherence rose from 0.682 before the guideline to 0.703 after
($P<.001$; Figure~\ref{fig:adherence-trend}), with monthly scores trending
upward across the period. The absolute change is small, and its significance
reflects the number of admissions rather than the size of the shift.
Component-level changes were uneven (Table~\ref{tab:adherence}). Glucose
monitoring improved most---on admission from 0.78 to 0.93---and discontinuation
of high-hypoglycemia-risk antidiabetic agents on admission rose from 0.62 to
0.75. Gains in the insulin-management components were significant but marginal
($+$0.01 to $+$0.03), leaving them close to their already low baselines, while
dose reduction after hypoglycemia did not change and SGLT2 discontinuation was
already complete in both periods. Center-level trajectories ranged from
improvement through no measurable change to modest decreases, so the pooled
increase conceals substantial inter-hospital heterogeneity.

\subsection{Adherence and clinical outcomes}

Higher adherence was associated with lower 30-day mortality
(Figure~\ref{fig:adherence-mortality}). Survivors had a higher mean adherence
score than those who died (0.693 vs 0.652, $P<.001$), and mortality fell across
adherence strata from 11\% to roughly 3\%. On the continuous score, each
0.1-point increase carried 16.8\% lower odds of death (OR, 0.832; 95\% CI,
0.811--0.853), so the association is graded rather than threshold-like, and
adjusting for age, sex, comorbidity burden and hospital left it intact
(adjusted OR, 0.757; 95\% CI, 0.727--0.788, per 1-SD increase). Adherence to the insulin-management
recommendations showed negligible correlation with subsequent hypoglycemic
burden, for both hypoglycemia ($r=0.052$) and severe hypoglycemia ($r=0.056$;
$n=40{,}719$).

\subsection{Outcome-associated recommendations}

Between 9 and 34 predictors were retained per outcome, but the number
surviving all three window definitions ranged from none to 32, tracking event
frequency: mortality yielded the most robust predictors and the rarer
complications few or none (Supplementary Table~S2). For mortality, the model
recovered well-established clinical predictors---age, red cell distribution
width, body mass index (BMI), dementia, cardiovascular comorbidity---which
gives the approach face validity, and several adherence components ranked among
the strongest retained predictors (Figure~\ref{fig:death-importance}). Across
all outcome models (Figure~\ref{fig:recurring-guidelines}), insulin following
hyperglycemia and admission glucose monitoring recurred most broadly, while
basal-dose components appeared in fewer, more outcome-specific models.
Direction was not uniform---the same component could carry opposite signs for
different endpoints.

\subsection{Propensity-weighted evaluation}

The overall-adherence analyses achieved satisfactory covariate balance (maximum
post-weighting SMD $<$0.08). Relative to the low-adherence group, above-median
adherence carried 44\% lower odds of 30-day mortality (OR, 0.557; 95\% CI,
0.514--0.603), 82\% lower odds of hyperglycemia (OR, 0.175) and 52\% lower odds
of severe hyperglycemia (OR, 0.485), all adjusted $P<.001$, against 20\% higher
odds of hypoglycemia (OR, 1.204; 95\% CI, 1.090--1.331;
Table~\ref{tab:iptw-overall}). Nothing survived correction for the remaining
five outcomes. The asymmetry between those two directions is the substantive
result: the reduction in hyperglycemic burden is roughly ninefold larger on the
log-odds scale than the accompanying increase in hypoglycemia, and the increase
in \emph{severe} hypoglycemia did not reach significance (OR, 1.129; adjusted
$P=.096$).

\begin{table}[t]
\centering
\caption{Inverse probability of treatment weighting analysis of overall
guideline adherence. Adherence was dichotomized at the median of each
outcome's eligible population, so $\beta$ is the log-odds ratio comparing the
above-median with the at-or-below-median group and OR $=\exp(\beta)$.
Confidence intervals are unadjusted; $P$ values were adjusted across outcomes
using the Benjamini--Hochberg false discovery rate procedure. Balance was
assessed by the maximum standardized mean difference (SMD) after weighting.
AKI, acute kidney injury; MI, myocardial infarction.}
\label{tab:iptw-overall}
\small
\setlength{\tabcolsep}{4pt}
\begin{tabular}{lrrlrr}
\toprule
Outcome & $\beta$ & OR & 95\% CI (OR) & Max SMD & $P$ \\
\midrule
Hyperglycemia        & $-1.742$ & 0.175 & 0.154 to 0.199 & 0.019 & $<.001$ \\
Severe hyperglycemia & $-0.724$ & 0.485 & 0.450 to 0.523 & 0.012 & $<.001$ \\
30-day mortality     & $-0.586$ & 0.557 & 0.514 to 0.603 & 0.020 & $<.001$ \\
Hypoglycemia         & $0.186$  & 1.204 & 1.090 to 1.331 & 0.034 & $<.001$ \\
Severe hypoglycemia  & $0.121$  & 1.129 & 1.001 to 1.274 & 0.030 & .096 \\
Hyperosmolarity      & $0.106$  & 1.112 & 0.652 to 1.896 & 0.020 & .800 \\
Ketoacidosis         & $0.116$  & 1.123 & 0.561 to 2.246 & 0.021 & .800 \\
Inpatient MI         & $0.055$  & 1.057 & 0.732 to 1.523 & 0.020 & .800 \\
Infection            & $-0.025$ & 0.975 & 0.809 to 1.175 & 0.020 & .800 \\
Inpatient AKI        & $0.018$  & 1.018 & 0.888 to 1.168 & 0.022 & .800 \\
\bottomrule
\end{tabular}
\end{table}

Only a subset of the recurrent component-level predictors remained significant
after weighting, indicating that several predictive associations were largely
explained by differences in baseline patient characteristics. A further six
significant associations were set aside because weighting failed to achieve
acceptable balance (maximum SMD, 0.90--2.46). All six concerned BMI or
creatinine measurement on admission, for which adherence is near-universal, so
that dichotomizing at the median produces groups too dissimilar to balance;
those coefficients were numerically degenerate and are reported as
non-estimable.

Eight associations were both significant and adequately balanced
(Table~\ref{tab:iptw-patterns}). Above-median adherence to routine basal
insulin dosing, routine glucose monitoring and insulin administration following
hyperglycemia each remained associated with lower 30-day mortality. For
glycemic outcomes, insulin following hyperglycemia (78\% lower odds; OR, 0.220)
and routine basal insulin dosing (60\% lower odds; OR, 0.402) were associated
with substantially lower odds of subsequent severe hyperglycemia, while insulin
following hyperglycemia also carried 28\% higher odds of hypoglycemia (OR,
1.283)---the same trade-off seen at the level of overall adherence. Adherence to glucose monitoring on
admission was associated with a higher recorded incidence of hyperglycemia and
severe hyperglycemia; since these outcomes are ascertained from glucose
measurements, more frequent early measurement mechanically increases the chance
of recording an event, so both associations likely reflect ascertainment rather
than harm.

\begin{table}[t]
\centering
\caption{Inverse probability of treatment weighting analysis of individual
guideline recommendations. Rows are the component--outcome pairs that were
both statistically significant after false discovery rate correction and
adequately balanced after weighting (maximum standardized mean difference
$<$0.10). A further six significant associations, all concerning BMI or
creatinine measurement on admission, failed the balance criterion (maximum
standardized mean difference, 0.90--2.46) and are not reported.
$^{\dagger}$Likely to reflect ascertainment rather than harm: the outcome is
detected from glucose measurements, so more frequent early measurement
increases the opportunity to record an event.}
\label{tab:iptw-patterns}
\small
\setlength{\tabcolsep}{4pt}
\begin{tabular}{llrrlr}
\toprule
Outcome & Recommendation & $\beta$ & OR & 95\% CI (OR) & $P$ \\
\midrule
30-day mortality & Routine basal dosage
  & $-0.291$ & 0.748 & 0.691 to 0.809 & $<.001$ \\
30-day mortality & Routine glucose monitoring
  & $-0.139$ & 0.870 & 0.809 to 0.937 & $<.001$ \\
30-day mortality & Insulin after hyperglycemia
  & $-0.098$ & 0.907 & 0.842 to 0.977 & .011 \\
\midrule
Severe hyperglycemia & Insulin after hyperglycemia
  & $-1.514$ & 0.220 & 0.202 to 0.239 & $<.001$ \\
Severe hyperglycemia & Routine basal dosage
  & $-0.912$ & 0.402 & 0.372 to 0.434 & $<.001$ \\
Severe hyperglycemia & Admission glucose monitoring$^{\dagger}$
  & $0.752$ & 2.121 & 1.974 to 2.280 & $<.001$ \\
\midrule
Hyperglycemia & Admission glucose monitoring$^{\dagger}$
  & $0.980$ & 2.664 & 2.389 to 2.971 & $<.001$ \\
\midrule
Hypoglycemia & Insulin after hyperglycemia
  & $0.249$ & 1.283 & 1.163 to 1.416 & $<.001$ \\
\bottomrule
\end{tabular}
\end{table}

These are the same components that showed the largest implementation gaps:
insulin following hyperglycemia (mean adherence, 0.15) and routine basal
insulin dosing (0.21--0.23) are simultaneously the least-implemented
recommendations assessed and, after weighting, among those most strongly
associated with mortality and with severe hyperglycemia.

% =====================================================================
% Table 2 -- merged overall adherence + pre/post change
% =====================================================================
% =====================================================================
% Table 3 -- IPTW, overall adherence
% =====================================================================
% =====================================================================
% Table 4 -- IPTW, individual recommendations
% =====================================================================
